\scanpage{38}
\mspage{37}

de la douzaine d'axiomes $\deleted{\unsure{p}}\mathrm{At\,1}-7$, $\mathrm{At\,L\,1}-3$, $\mathrm{At\,\deleted{\unsure{L}}CL(spac)}$.

Si on se donne $\parallel$ sur $M$, et non sur $\unsure{L|\mathcal L}$, il faut simplement rajouter un axiome supplémentaire :

$(\mathrm{ML}\,\deleted{\unsure{5}}6)$ Soit $\unsure{X|x},Y\in\deleted{\unsure{L}}M$. Alors
\[
 X\parallel Y\quad\Longleftrightarrow\quad \forall\,x\unsure{\in|,}y\deleted{\unsure{-}}\in L,\quad
 x\ll X,\ y\deleted{\unsure{\leq}}\ll Y\ \Rightarrow\ x\parallel y.
\]

\uline{Variante pour atlas polyédraux}. \inserted{(satisfaisant $\mathrm{At\,1}-\mathrm{At\,7}$, $\mathrm{At\,L\,1}-\mathrm{At\,L\,3}$, \inserted{$\mathrm{At\,CL(pol)}$})} C'est exactement la \unsure{$=$} chose, sauf qu'on change en conséquence la définition de $X\mathrel{\unsure{\ooalign{\hfil$>$\hfil\cr\hfil$<$\hfil\cr}|\gtrless}}Y$ (p.~29), et qu'il faut supposer sur $M,\leq$ une condition supplémentaire

$(\mathrm{ML\,0\ pol})$ $\deleted{\unsure{\exists}}\forall\,X\in M$, $\widetilde X=M_{\leq X}$ est un

\noindent\begin{tabular}{@{}|p{\dimexpr\linewidth-2\tabcolsep-1pt\relax}@{}}
ens ordonné \uline{polyédral} (p.~\unsure{13|15}), \unsure{i.e|ie} $\deleted{\exists\,\unsure{X|x},Y,Z}$ deux \unsure{éléments|élmts} $Y,Z$ de $\deleted{\unsure{X}}M$, \inserted{$\leq$} qui sont majorés (par un $X$) et minorés\unsure{,} \deleted{\unsure{adme}} (par un $T$) admettent un inf -- ou \unsure{encore}, les ensembles $M_{T\leq\unsure{\cdot|\bullet}\leq\unsure{x|X}}$ ont des inf finis (donc aussi des \unsure{sup|type} finis s'ils sont finis).
\end{tabular}
